%%%%%%%%%%%%%%
% Modelo XXXXXXXXXXXXXXX
% Autor: Leonardo Garcez Dalenogari Alba
% Ano: 2022
% Contato: leonardo.2017007421@aluno.iffar.edu.br

% IMPORTANTE: 
%   O modelo utiliza o compilador XeLaTex (talvez possam ser utilizados outros, mas não é garantido o correto funcionamento, visto que seria necessário testar, porém, eu sei que alguns vão apresentar problemas na hora de compilar, como o pdflatex)
%   Se você estiver utilizando o Overleaf, com o documento aberto, clique em Menu, no canto superior esquerdo, e selecione o XeLaTex como compilador
%   Também é utilizado o biber para gerenciar as referências bibliográficas
%%%%%%%%%%%%%%

\documentclass[
    12pt, % Tamanho de fonte definido para 12, se você precisar utilizar um texto em tamanho 10, utilize o comando \footnotesize, que será impresso em tamanho 10
    oneside,
    %twoside, % A opção permite gerar o documento utilizando o verso das páginas, porém, ela não indica à impressora que você está utilizando isso ou que ela deva imprimir nos dois lados das folhas (fonte: https://en.wikibooks.org/wiki/LaTeX/Document_Structure)
    %openany,
    %openright,
    a4paper,
    english,
    brazil
]{report}

% Pacotes que envolvem codificação do conteúdo. Essenciais.
\usepackage[T1]{fontenc}
% \usepackage[utf8]{inputenc} % Não é necessário se você estiver utilizando um compilador LaTeX baseado já em UTF-8

% Responsável pela quebra silábica de palavras no final de linha de forma correta para o português
\usepackage[brazil]{babel}
% O pacote babel recomenda o uso do csquotes, para garantir que textos citados estejam escritos corretamente de acordo com as regras do seu idioma.
\usepackage[portuguese=brazilian]{csquotes}

%% Pacotes para formatação de fonte do documento
\usepackage{fontspec}
% Pode ser Times New Roman:
% \setmainfont{Times New Roman} 
% Ou pode ser Arial, também:
\setmainfont{Arial}

% Pacote responsável por adicionar símbolos matemáticos. Símbolos, por exemplo, também necessários para o correto referenciamento de leis e decretos
\usepackage{amssymb}
% Só que também é importante realizar a configuração para adicionar os comandos que imprimem esses símbolos associados aos caracteres que imprimem, assim, você pode adicionar o símbolo diretamente no texto
\usepackage{newunicodechar} % Pacote que permite definir novos caracteres
\usepackage{config/unicode-char} % Arquivo de configuração com a listagem desses caracteres

%% Pacotes para formatação de conteúdo na página
\usepackage[left=3cm,top=3cm,right=2cm,bottom=2cm]{geometry} % Pacote responsável por formatar as dimensões das margens das páginas.
\usepackage{setspace} % Adiciona opções de espaçamento de texto
\usepackage{ragged2e} % Adiciona opções mais completas para alinhamento de texto
\usepackage{indentfirst} % Faz com que o primeiro parágrafo de cada seção possua indentação também

% Pacote que permite criar estilos para o cabeçalho
\usepackage{fancyhdr}
\usepackage{config/cabecalho} % Cabeçalho customizado para o modelo

%% Adiciono os pacotes para configurar e estilizar o sumário e as seções
% \usepackage{tocloft}
\usepackage{titletoc} % Não consegui criar pelo titletoc a parte de todas as seções estarem alinhadas, então vou usar o tocloft
\usepackage{titlesec}
\usepackage{alphalph} % Permite lidar com contadores que possuam valor maior que 26 utilizando o alfabeto (o número de letras), ou seja, se passar de 26, vai utilizar AA, AB, AC... Nativamente o LaTeX não realiza esse controle
\usepackage{config/apendice-anexo} % Criação dos comandos para apêndice e anexo

%% Agora customizo as seções e sumário para seguirem as normas do IFFar (itens adicionais do sumário, como apêndice e anexos, também são configurados nesses arquivos)
\usepackage{config/sumario} % Importa o sumário feito para o modelo
%% Adiciona a possibilidade de estilizar os títulos de seções, subseções e afins
\usepackage{config/titulo-sections} % Importa as estilizações feitas das seções e subseções para o modelo

% Pacote para facilitar a customização de listas numeradas, que serão utilizadas para as alíneas e subalíneas
\usepackage{enumitem}
\usepackage{config/enumerate} % Definição das listas de alíneas, subalíneas e índice

\usepackage{float} % Adiciona configurações adicionais para elementos float (figuras, tabelas, quadros, etc.)
\usepackage{tabularx} % Tabelas no LaTeX são limitadas em algumas características, este pacote resolve algumas dessas falhas
\usepackage{graphicx} % Adiciona a capacidde de adicionar figuras, além de outras funcionalidade que podem ser utilizadas em quadros e tabelas.
% Figuras podem ser imagens (SVG, por padrão, não funciona), mas documentos PDF também são compatíveis
\graphicspath{ {./img/} } % Adiciona um caminho padrão para arquivo
% Permite modificar a apresentação das legendas (caption) de figuras, quadros e tabelas
\usepackage[labelsep=endash, font=singlespacing, justification=centering]{caption} % Modifica a separação da numeração da legenda com a descrição dela
\usepackage{newfloat} % Permite criar floats, como o caso dos quadros, de forma prática
\DeclareFloatingEnvironment[listname={Lista de quadros}, name={Quadro}]{quadro} % Crio o elemento float Quadro
% Removo a contabilização de capítulo na numeração das figuras, tabelas e quadros (isso ocorre ao utilizar o documentclass de book ou report)
\counterwithout{figure}{chapter} 
\counterwithout{table}{chapter}
\counterwithout{quadro}{chapter}
% Obs.: Se você precisar criar outro elemento (e outra lista), como de gráficos ou de qualquer outro tipo, é, teoricamente, bastante simples: 
%   - Você declara o novo ambiente com o \DeclareFloatingEnvironment, junto com o nome do elemento e o título de sua lista
%   - Remove a contabilização do capítulo com o \counterwithout (apesar que acho que isso pode ser realizado no próprio \DeclareFloatingEnvironment)
%   - E depois vai no final do arquivo config/sumario.sty e copia a estilização usada para os elementos de figura ou tabela, realizando a modificação do nome do elemento

%% Pacotes para adicionar o uso de citações e referências no padrão ABNT
% Permite customizar as notas de rodapé. Necessário para a indentação das notas
\usepackage[hang]{footmisc}
\usepackage{config/rodape} % Configurações das notas de rodapé
% Usa o biber por ser uma versão mais parruda de biblatex, que é uma versão mais parruda do bibtex, ou algo assim. É importante por resolver uma quantia de erros por conta de caracteres, principalmente na URL das referências
\usepackage[backend=biber, style=abnt]{biblatex}
\addbibresource{bibliografia.bib} % Você pode adicionar mais de um arquivo.bib, apenas adicione outra linha com o nome do segundo arquivo
\usepackage{config/bibliografia-abnt} % Ajustes necessários por causa que o pacote biblatex-abnt não foi atualizado para algumas das mais novas normas da ABNT

% Pacote responsável por permitir criar o índice. Se você não for utilizar índice, e provavelmente não vai, é só comentar essas duas linhas abaixo e o comando \printindex no final do documento
\usepackage{imakeidx}
\makeindex[columns=3, title=Índice]

%Pacote responsável por permitir colocar url. Ele cria link de atalho no sumário, citações e referência de figuras, tabelas, etc.
\usepackage[hidelinks]{hyperref}
\hypersetup{
    unicode=true
}
\usepackage{url}

\usepackage{lipsum} % Pacote meramente cosmético, para testes. Remova ele quando for utilizar o modelo.

\begin{document}
\onehalfspacing
\justifying
\setlength\parindent{1.25cm}% Configura o espaçamento da indentação dos parágrafos
\setlength\parskip{1.241pt}% Configura o espaço separação entre parágrafos.
% Dados de identificação do trabalho
\newcommand{\tituloTrabalho}{Modelo LaTeX}
\newcommand{\separadorTrabalho}{:\space} % Aqui você identifica o caractere que irá separar o título do subtítulo. Eu já vi trabalho que não utiliza dois pontos, então sei lá, deixei para permitir configuração, recomendo conversar com quem te orienta. Se não possuir subtítulo, deixar em branco
\newcommand{\subtituloTrabalho}{como é gostoso usar um template redondinho} % Se não possuir subtítulo, deixar em branco

% O seu nome completo
\newcommand{\autor}{Insira o seu nome completo}

% Os dados sobre a apresentação da defesa
\newcommand{\cidadeDefesa}{Insira a cidade onde ocorre a defesa}
\newcommand{\anoEntrega}{2023} % O ano em que o trabalho foi entregue
\newcommand{\diaDefesa}{1}
\newcommand{\mesDefesa}{janeiro}
\newcommand{\anoDefesa}{2023} % O ano em que o trabalho foi defendido (atente-se que pode ter diferença)
\newcommand{\tituloCurso}{Insira o título de formação do seu curso}
\newcommand{\nomeCurso}{Insira o título de formação do seu curso}

% Os dados de quem te orientou
\newcommand{\orientador}{Insira o nome do(a) orientador(a)}
\newcommand{\orientadorTitulacao}{Prof.(a) Dr.(a)}
\newcommand{\coorientador}{Insira o nome do(a) coorientador(a)} % Deixe em branco se não possuir
\newcommand{\coorientadorTitulacao}{Prof.(a) Dr.(a)}

% Os dados sobre a sua banca
\newcommand{\banca1}{Insira o nome da primeira pessoa da banca}
\newcommand{\banca1Titulacao}{Prof.(a) Dr.(a)}
\newcommand{\banca2}{Insira o nome da segunda pessoa da banca}
\newcommand{\banca2Titulacao}{Prof.(a) Dr.(a)}
\newcommand{\banca3}{Insira o nome da terceira pessoa da banca} % Deixe em branco a terceira pessoa da banca se for um TCC (ou tanto faz, não vai ser impresso de qualquer forma) 
\newcommand{\banca3Titulacao}{Prof.(a) Dr.(a)}

MINISTÉRIO DA EDUCAÇÃO
INSTITUTO FEDERAL DE EDUCAÇÃO, CIÊNCIA E TECNOLOGIA
FARROUPILHA

\begin{Center}
% Imprimir a marca do IFFar caso seja tese ou dissertação
\tccOuNao{}{
    \begin{figure}[H]
        \Centering\singlespacing
        \includegraphics[width=3cm]{marca-iffar-vertical.png}
    \end{figure}
}

% MINISTÉRIO DA EDUCAÇÃO\\
\MakeUppercase{\nomeInstituicao}

\MakeUppercase{\nomeCampus}

\MakeUppercase{\nomeCursoPrograma}

\vspace{\fill}

\autorTrabalho

\vspace{\fill}

\textbf{\tituloTrabalho\separadorTrabalho}\subtituloTrabalho

\vspace{\fill}

\cidadeDefesa

\anoEntrega
\end{Center}
\OnesideTwoside{\clearpage}{\cleardoublepage} % ITEM OBRIGATÓRIO
\begin{Center}
\autorTrabalho

\vspace*{\fill}

\textbf{\tituloTrabalho\separadorTrabalho}\subtituloTrabalho
\end{Center}

\vspace*{\fill}

% Tamanho total da página: 21cm
% Margem esquerda: 3cm
% Margem direita: 2cm 
% Espaço para ocupar a metade restante da página: 8cm
\hspace{\fill}
\begin{minipage}{8cm}
    \footnotesize
    \begin{singlespace}
        \tccOuNao{ 
            % Natureza do trabalho para TCC
            Trabalho de Conclusão de Curso de Graduação em \modalidadeCurso\space em \nomeCurso do \nomeCampus\space do \nomeInstituicao\space como requisito para a obtenção do título de \tituloCurso\space em \nomeCurso.\\
            Orientador(a) \orientadorTitulacao\space \orientador
            \isEmpty{\coorientador}{\\Coorientador(a) \coorientadorTitulacao \coorientador}
        }{ 
            % Natureza do trabalho para tese ou dissertação
            \tipoTrabalho\space submetida ao Programa de \nomeCurso\space do Instituto Federal Farroupilha para a obtenção do título de \tituloCurso\space em [.............]\\
            Orientador(a) \orientadorTitulacao\space \orientador
            \isEmpty{\coorientador}{\\Coorientador(a) \coorientadorTitulacao \coorientador}
        }
    \end{singlespace}
\end{minipage}

\vspace*{\fill}

\begin{Center}
\cidadeDefesa

\anoEntrega
\end{Center}  % ITEM OBRIGATÓRIO
\tccOuNao{ % Folha de aprovação, para o caso de TCC
    \begin{Center}
        \autorTrabalho

        \vspace{\baselineskip}
        \textbf{\tituloTrabalho\separadorTrabalho}\subtituloTrabalho
        
        \vspace{\baselineskip}
        Este Trabalho de Conclusão de Curso foi julgado adequado para obtenção do Título de \tccTituloCurso\space e aprovado em sua forma final pelo Curso \tccNomeCursoCompleto.

        \vspace{\baselineskip}
        \cidadeDefesa, \diaDefesa\space de \mesDefesa\space de \anoDefesa.

        \vspace{\tamanhoAssinatura}
        \rule{6.5cm}{1pt}\\ 
        \coordenadorTitulacao\space \coordenador\\ 
        \coordenadorApresentacao\space do Curso

        \vspace{\baselineskip}
        \textbf{Banca examinadora:}

        \vspace{\tamanhoAssinatura}
        \rule{6.5cm}{1pt}\\ 
        \orientadorTitulacao\space \orientador\\ 
        \orientadorApresentacao\\ 
        \orientadorInstituicao

        % Verifico se o coorientador está preenchido, para imprimir
        \notEmpty{\coorientador}{
            \vspace{\tamanhoAssinatura}
            \rule{6.5cm}{1pt}\\ 
            \coorientadorTitulacao\space \coorientador\\ 
            \coordenadorApresentacao\\ 
            \coorientadorInstituicao
        }

        \vspace{\tamanhoAssinatura}
        \rule{6.5cm}{1pt}\\ 
        \bancaUmTitulacao\space \bancaUm\\ 
        \bancaUmApresentacao\\ 
        \bancaUmInstituicao

        \vspace{\tamanhoAssinatura}
        \rule{6.5cm}{1pt}\\ 
        \bancaDoisTitulacao\space \bancaDois\\ 
        \bancaDoisApresentacao\\ 
        \bancaDoisInstituicao
    \end{Center}
}{ % Folha de certificação, para o caso de teses ou dissertações
    \begin{Center}
        \autorTrabalho

        \textbf{\tituloTrabalho\separadorTrabalho}\subtituloTrabalho
    \end{Center}
    
    \vspace{\baselineskip}
    O presente Trabalho de \teseTipoFormacao\space foi avaliado e aprovado por banca examinadora composta pelos seguintes membros:

    \begin{Center}
        \vspace{\baselineskip}
        % \rule{6.5cm}{1pt}\\
        \orientadorTitulacao\space \orientador\space - \MakeLowercase{\orientadorApresentacao}\\  
        \orientadorInstituicao

        % Verifico se o coorientador está preenchido, para imprimir
        \notEmpty{\coorientador}{
            \vspace{\baselineskip}
            % \rule{6.5cm}{1pt}\\
            \coorientadorTitulacao\space \coorientador\space - \MakeLowercase{\coorientadorApresentacao}\\ 
            \coorientadorInstituicao
        }

        \vspace{\baselineskip}
        % \rule{6.5cm}{1pt}\\
        \bancaUmTitulacao\space \bancaUm\space - \MakeLowercase{\bancaUmApresentacao}\\ 
        \bancaUmInstituicao

        \vspace{\baselineskip}
        % \rule{6.5cm}{1pt}\\
        \bancaDoisTitulacao\space \bancaDois\space - \MakeLowercase{\bancaDoisApresentacao}\\    
        \bancaDoisInstituicao

        \vspace{\baselineskip}
        % \rule{6.5cm}{1pt}\\
        \bancaTresTitulacao\space \bancaTres\space - \MakeLowercase{\bancaTresApresentacao}\\    
        \bancaTresInstituicao
    \end{Center}

    \vspace{\baselineskip}

    Certificamos que esta é a \textbf{versão original e final} do trabalho de conclusão que foi julgado adequado para obtenção do título de \teseTitulo\space em \teseFormacao.
    
    \begin{Center}
        \vspace{\tamanhoAssinatura}
        \rule{6.5cm}{1pt}\\
        \coordenadorTitulacao\space \coordenador\\
        \coordenadorApresentacao\space do Programa de Pós-Graduação

        \vspace{\tamanhoAssinatura}
        \rule{6.5cm}{1pt}\\
        \orientadorTitulacao\space \orientador\\
        \orientadorApresentacao

        \vspace{\fill}
        \cidadeDefesa, \anoEntrega.
    \end{Center}
}

% Recomendação do Guia de Normalização de iniciar cada elemento pré-textual e textual no anverso da página
\OnesideTwoside{\clearpage}{\cleardoublepage}  % ITEM OBRIGATÓRIO
\begin{SecoesNaoNumeradas}
  \vspace*{\fill}

% Tamanho total da página: 21cm
% Margem esquerda: 3cm
% Margem direita: 2cm
% Recuo necessário: 4cm
% Espaço restante: 12cm
\hspace{\fill}
\begin{minipage}{12cm}
    \setlength{\parindent}{1.25cm}
    
    Este trabalho é dedicado [especificar o texto da dedicatória, escrito pelo(a) autor(a), registrando a quem dedica o Trabalho Acadêmico-Científico realizado ou prestando uma homenagem a alguém que considere importante no transcorrer do seu processo formativo (pessoal e/ou acadêmico-profissional), ainda que essa pessoa não tenha contribuído diretamente para os resultados obtidos no trabalho].
\end{minipage}

% Recomendação do Guia de Normalização de iniciar cada elemento pré-textual e textual no anverso da página
\OnesideTwoside{\clearpage}{\cleardoublepage}

  \chapter*{Agradecimentos}
\thispagestyle{empty}

\lipsum[2]
  \vspace*{\fill}

% Largura da folha: 21cm
% Margem à esquerda: 3cm
% Margem à direita: 2cm
% Espaço para texto: 16cm;
% Margem à esquerda necessária: 10cm
% Se for margem de 10cm da página inteira, tamanho de: 11cm
% Se for margem de 10cm da parte de texto, sem considerar os 3cm: 6cm 
\hspace{\fill}
\begin{minipage}{6cm}
    \begin{singlespace} % A epígrafe tem que usar espaçamento símples
        % Seu texto recomenda-se em itálico e autor em negrito
        \textit{Lorem ipsum dolor sit amet, consectetur adipiscing elit. Curabitur et neque metus. Ut volutpat sem quis lacus rutrum, a viverra odio tincidunt. Quisque eu mollis enim. Suspendisse porta vehicula pharetra. Sed magna risus, consequat vitae tempus at, pulvinar ut sapien. Donec sit amet metus feugiat massa molestie venenatis id vel est. Fusce dapibus nisi vitae urna suscipit, non hendrerit mauris consequat. Nunc in quam vel dolor auctor interdum eget nec metus. Donec ornare libero euismod mi vehicula tincidunt. Pellentesque fermentum vulputate lacus, sit amet facilisis dolor faucibus et} (\textbf{Autor}).
    \end{singlespace}
\end{minipage}

% Recomendação do Guia de Normalização de iniciar cada elemento pré-textual e textual no anverso da página
\OnesideTwoside{\clearpage}{\cleardoublepage}
  \chapter*{Resumo}\myIdxResumoVernac

% O \noindent é necessário por causa que, como eu configurei para indentar o primeiro parágrafo no texto, como exigido, o LaTeX iria normalmente indentar o parágrafo
% Só escrever o teu resuminho, sem estresse
\noindent [Inserir o resumo que representa uma síntese do trabalho e cujo texto é escrito na língua vernácula, de modo objetivo, sucinto e claro. Deve ser digitado como um bloco único de parágrafo e apresentar uma sequência de frases concisas, esclarecendo: a delimitação da temática/assunto abordado; o objetivo principal da pesquisa; a questão norteadora ou hipótese(s) do trabalho, bem como o viés teórico-metodológico que o fundamentou. Deve constar também o método e os procedimentos técnicos utilizados na obtenção e na análise dos dados; apontando principais resultados e pontos de discussão. O resumo não deve ser uma enumeração de tópicos, nem conter citações. Recomenda-se isar o verbo na voz ativa e na terceira pessoa do singular. O uso do verbo na primeira pessoa depende do diálogo com o(a) orientador(a) e o Programa de Pós-graudação no qual o Trabalho Acadêmico-Científico está sendo desenvolvido. \textbf{Seu texto deve ser digitado com espaçamento entre linhas simples}\footnote{O próprio Guia de Normalização contradiz isso. E a NBR 14724 não cita o resumo como um dos casos de espaçamento simples.} e fonte de tamanho 12. Evitar o uso de fórmulas, diagramas e siglas sem a respectiva designação. O resumo deve conter de 150 a 500 palavras e ser elaborado em conformidade com a NBR 6028:2021. Abaixo do resumo, deve-se deixar uma linha em branco e colocar a expressão ``Palavras-chave'', seguida de dois pontos e um espaço. Em seguida, deve-se informar as palavras-chave, especificando de três a cinco delas (constituída por termos extraídos de Thesaurus da área ou expressões significativas ao texto, preferencialmente que não sejam coincidentes com as palavras usadas no título do trabalho para favorecer sua busca por mecanismos de recuperação). Exceto quando ela se constituir como um substantivo próprio ou nome científico, cada palavra-chave é iniciada por letra minúscula, seguida de ponto e vírgula, e de um espaço, sendo a última finalizada por um ponto]. 

\vspace{\baselineskip} % Fazer o espacinho entre as duas partes

\noindent \textbf{Palavras-chave:} palavra-chave 1; palavra-chave 2; palavra-chave 3.

% Recomendação do Guia de Normalização de iniciar cada elemento pré-textual e textual no anverso da página
\OnesideTwoside{\clearpage}{\cleardoublepage} % ITEM OBRIGATÓRIO
  % É exatamente o mesmo código do 
% Só fico na dúvida, será que o abstract tem que ser escrito em itálico, já que é todo em outro idioma? Se sim, é só englobar todas as partes com texto (sem o \noindent) em um \textit{}
\chapter*{Abstract}


% O \noindent é necessário por causa que, como eu configurei para indentar o primeiro parágrafo no texto, como exigido, o LaTeX iria normalmente indentar o parágrafo
\noindent \lipsum[1] % Só escrever o teu abstractzinho, sem estresse

\vspace{\baselineskip} % Fazer o espacinho entre as duas partes

\noindent \textbf{Keywords:} keyword 1; keyword 2; keyword 3.

% Recomendação do Guia de Normalização de iniciar cada elemento pré-textual e textual no anverso da página
\OnesideTwoside{\clearpage}{\cleardoublepage} % ITEM OBRIGATÓRIO

  \renewcommand\addvspace[1]{} % Remove o espaço adicional nos títulos das listas de figuras, tabelas e quadros
  % Lista de figuras:
  \listoffigures
  % Lista de tabelas:
  \listoftables
  % Lista de quadros:
  \listofquadros

  % Crio um elemento utilizado para listar cada abreviatura e sigla
% Obs.: Dependendo do tamanho das siglas utilizadas, você pode querer aumentar ou diminuir o espaço entre a sigla e o seu significado, porém, mantendo o alinhamento, para isto, você modifica o comando \newcommand\siglaLargura{10ex}, substituindo o 10ex pelo tamanho que você deseja
% O tamanho 10ex suporta uma sigla de até mais ou menos 10 letras, visto que a unidade utiliza como referência o tamanho da letra "x". Você também pode utilizar a unidade "em", que utiliza como referência a letra "m", se quiser.
% Fonte: https://tex.stackexchange.com/a/508154
\newcommand\siglaLargura{10ex} % Comprimento do espaço onde fica a sigla
\newcommand\siglaGap{1ex} % Comprimento do vão entre sigla e seu significado, para dar uma folguinha no tamanho máximo
\newcommand\nomeSiglaLargura{\dimexpr\linewidth-\siglaLargura-\siglaGap\relax}
\newcommand\sigla[2]{\noindent\parbox[t]{\siglaLargura}{#1\strut}%
  \hspace{\siglaGap}%
  \parbox[t]{\nomeSiglaLargura}{#2\strut}}

% Adicione as suas siglas:
\chapter*{Lista de abreviaturas e siglas}
% Para adicionar uma abreviatura ou sigla na lista apenas precisa adicionar a sigla dentro do \sigla{<insira a sigla>}{<insira o nome completo da sigla>}

\sigla{IBGE}{Instituto Brasileiro de Geografia e Estatística}

\sigla{IFFar}{Instituto Federal de Educação, Ciência e Tecnologia Farroupilha}

\sigla{TCC}{Trabalho de Conclusão de Curso}

\sigla{T}{Teste}

\sigla{xxxxxxxxxx}{Lorem ipsum dolor sit amet, consectetur adipiscing elit. Curabitur et neque metus. Ut volutpat sem quis lacus rutrum, a viverra odio tincidunt}
  % Crio um elemento utilizado para listar cada símbolo e seu significado
% Obs.: É a mesma lógica da lista de siglas. Ajuste o tamanho do comprimento da "caixinha" do símbolo conforme a sua necessidade
% Fonte: https://tex.stackexchange.com/a/508154
\newcommand\simboloLargura{5ex} % Comprimento do espaço onde fica o símbolo
\newcommand\simboloGap{1ex} % Comprimento do vão entre símbolo e seu nome, para dar uma folguinha no tamanho máximo
\newcommand\simboloNomeLargura{\dimexpr\linewidth-\simboloLargura-\simboloGap\relax}
\newcommand\simbolo[2]{\noindent\parbox[t]{\simboloLargura}{#1\strut}%
  \hspace{\simboloGap}%
  \parbox[t]{\simboloNomeLargura}{#2\strut}}

% Adicione os seus símbolos:
\chapter*{Lista de símbolos}

% Para adicionar uma abreviatura ou sigla na lista apenas precisa adicionar o símbolo no comando \simbolo{<insira o símbolo>}{<insira o significado do símbolo>}

\simbolo{\LaTeX}{LaTeX}

\simbolo{\S}{Parágrafo}

\simbolo{$\Sigma$}{Somatória}

\simbolo{$\sum_{n = 1}^{\infty}$}{Matemágica}

\simbolo{xxxxx}{Lorem ipsum dolor sit amet, consectetur adipiscing elit. Curabitur et neque metus. Ut volutpat sem quis lacus rutrum, a viverra odio tincidunt}

% Recomendação do Guia de Normalização de iniciar cada elemento pré-textual e textual no anverso da página
\OnesideTwoside{\clearpage}{\cleardoublepage}

  %Sumário
  \tableofcontents  % ITEM OBRIGATÓRIO
\end{SecoesNaoNumeradas}
% \clearpage

\begingroup
  \chapter{Introdução}
Este é um modelo para trabalhos acadêmicos do Instituto Federal de Educação, Ciência e Tecnologia Farroupilha (IFFar) construído no \LaTeX{} e, assim como qualquer modelo, possui suas observações. Por trabalho acadêmico, refere-se aos Trabalhos de Conclusão de Curso (TCC), teses e dissertações e, sobre o modelo, por ser construído no \LaTeX, ele é fácilmente adaptável, permitindo estender suas implementações. Um outro aspecto desse modelo, também, é que, de início, é necessário realizar o preenchimento de alguns dados, para a configuração de alguns elementos pré-textuais, que será melhor explicado no segundo capítulo.
\endgroup

\begingroup
  \chapter{Configurações e discussões sobre o modelo}\label{capitulo:configuracao}
\section{Preenchimento de dados}
Uma das primeiras coisas que deve ser feita neste modelo é o preenchimento dos dados utilizados pelo modelo. Isso é realizado no arquivo \texttt{elementos/DADOS.tex} e, este arquivo, foi pensado como o ponto central para configuras algumas páginas pré-textuais, como capa e folha de rosto, pensando em facilitar a alteração desses registros. Se a pessoa quiser, pode realizar o preenchimento dos dados manualmente em cada um desses arquivos.

A pessoa deve iniciar indicando se o seu trabalho é do tipo TCC ou não, indicando no comando \verb|\isTCC| o valor \textbf{1 = TCC} ou \textbf{0 = Tese ou Dissertação}. Esse comando é necessário para o modelo alterar a capa e folha de rosto, além de alternar entre folha de aprovação, para TCCs, ou folha de certificação, para teses e dissertações. Após isso, podem ser preenchidos alguns dados gerais utilizados entre as páginas, seguido pela escolha do preenchimento de dados: para TCC; ou Tese e Dissertação. Por último, são preenchidos os dados sobre a sua banca e coordenação do seu curso/Programa.

Importante ressaltar de que, no caso de não existir coorientação, o comando \verb|\coorientador| seja deixado totalmente vazio. Pois o modelo verifica o preenchimento desse valor para adicionar ou remover suas citações nas páginas. Além disso, realize a alteração no tamanho do espaço vertical dado para os campos de assinatura\footnote{Apesar de você adicionar o determinado tamanho para os campos, ele não vai ter especificamente aquele tamanho utilizado. Isso acontece por causa que ainda são somados aos espaços entre parágrafos, além do espaço de uma linha do filete de assinatura.}. Isso é importante, pois, o espaço foi pensado no modelo utilizando a fonte Arial, de tamanho maior se comparado com a Times New Roman, e considerando a existência de coorientação no trabalho. E, por último, não tenha medo de realizar ajustes manualmente em algumas dessas páginas pré-textuais, caso a lógica utilizada com os comandos para o preenchimento automático não funcione tão bem.

\section{Ajustando outros espaçamentos}
Outros elementos que podem receber o ajuste em determinados espaçamentos são: sumário; lista de siglas; lista de símbolos; e a margem das notas de rodapé. Neles foram utilizados a unidade \textbf{ex}, que usa como referência o tamanho de uma letra \textit{x}, mas você pode utilizar qualquer outra medida aceita pelo \LaTeX{}. Isso é algo extremamente opcional de se realizar, porém, podem existir situações que necessitam disso, pois o espaço utilizado ou é muito grande, ou muito pequeno.

Para ajustar o espaçamento do \textbf{sumário}: Acesse o arquivo \\\texttt{config/sumario.sty} e altere o comando \verb|\espacoAlinhamentoSumario|. O \texttt{10ex} representa o espaço utilizado, que foi definido considerando um possível número de seções, que chegam até o quinto nível.

Para ajustar as listas de \textbf{símbolos} e de \textbf{siglas}: Acesso os respectivos arquivos de cada uma dessas duas listas, na pasta \texttt{elementos/}, alterando o espaço definido nos comandos \verb|\simboloLargura| e \verb|\siglaLargura|, respectivamente.

O ajuste do espaçamento do algarismo da nota de rodapé pode ser alterado no arquivo \texttt{config/rodape.sty}. Só o comando \verb|\footnotemargin| no arquivo e escolher o tamanho do espaçamento para alinhamento. Não foi identificado no Guia de Normalização e nem em normas da ABNT, ou então passou despercebido, a indicação de qual deveria ser a margem entre o algarismo e a nota, então, a configuração utilizou o número que melhor se enquadrou.

\section{Escolha de fonte}
No início do arquivo \texttt{main.tex} você pode optar entre utilizar fonte Arial ou Times New Roman. Isso é realizado pelo o comando \verb|\setmainfont|, você apenas precisa localizar esses comandos, comentar uma dessas linhas e descomentar a fonte escolhida no início do arquivo. Uma condição opcional é, no caso do uso da Arial, a utilização da Andale Mono como a opção de fonte monoespaçada, utilizada em algumas situações. Caso opte por utilizá-la e esteja utilizando o \LaTeX{} localmente em sua máquina, baixe e instale a fonte\footnote{Local por onde se pode baixá-la: \url{https://www.cufonfonts.com/font/andale-mono}}. Para o Overleaf nenhum trabalho adicional é necessário, pois, ele já apresenta suporte à fonte. Se não for utilizar a fonte monoespaçada, apenas comente a sua linha abaixo da definição da Arial.

\subsection{Configuração de fontes para matemática}
Para o caso de utilização de equações e fórmulas, ainda é necessária a seleção da fonte utilizada nesses elementos. No modelo, é utilizado o pacote \texttt{unicode-math} para a configuração de fontes e, dependendo da seleção anterior de fonte (Arial ou Times New Roman), também é necessário selecionar as suas respectivas fontes. Se você não for utilizar os elementos matemáticos do \LaTeX{}, pode ignorar essa seção e remover a importação\footnote{Recomendo que apenas comente as linhas} dos pacotes \texttt{amsmath} e \texttt{unicode-math}, além de suas definições de fontes matemáticas abaixo\footnote{O pacote \texttt{amssymb}, e seu arquivo de configuração abaixo, também pode ser removido, mas existe o seu uso na definição de alguns caracteres, para permitir a utilização do caractere diretamente no texto dos arquivos.}.

Para o caso de Times New Roman, é utilizada a Latin Modern Math como fonte matemática, com a substituição da Times New Roman para os números e letras latinas e gregas. Já, para o caso da utilização da Arial, ocorre o mesmo processo de substituição, mas utilizando a fonte matemática sem serifa Fira Math. Se estiver utilizando o \LaTeX{} localmente, assim como no caso da Andale Mono, será necessário baixar a fonte e instalá-la em seu computador\footnote{https://github.com/firamath/firamath/releases}.

Ao compilar, o \LaTeX{} irá enviar um conjunto de alertas, por usar a fonte Times New Roman ou Arial em alguns conjuntos como \textit{mathfont}. Não foi identificado problemas que ocorram por causa do uso dessas fontes, porém, se você identificar problemas em suas equações, você pode: no caso de uso da Times New Roman, comentar as linhas de \verb|\setmathfont| que utilizem a Times, pois a Latin Modern já é uma fonte semelhante o suficiente; e no caso da Arial, também comentar as linhas de \verb|\setmathfont| utilizando Arial, porém, a Fira Math pode não ser tão semelhante a ela, então, sugere-se a troca da fonte matemática\footnote{As fontes suportadas pelo Overleaf estão em \textcite{site:overleaf-fontes-2023}. E no link a seguir estão um conjunto de exemplos de fontes com suporte matemático: \textcite{site:tex-fontes-2022}}.

\section{Criação de novos elementos de ilustração e listas}\label{section:novos-elementos}
Uma característica que pode ser necessária em seu trabalho é a criação de novos elementos de ilustração para listagem. Por exemplo, você pode necessitar criar um elemento \iindex{Gráfico}{Ilustracao@Ilustração!Grafico@}, ou Fotografias, junto com suas respectivas listas. E isso é facilmente realizavel no arquivo \texttt{main.tex}: abaixo da importação do pacote \texttt{newfloat} estão contidas as instruções para realizar tal ajuste.

\section{Utilização do \LaTeX{} no PC ou pelo navegador}
Para utilizar o \LaTeX{} você basicamente tem duas opções: instalar na sua máquina; ou utilizar pelo navegador. Utilizando pelo navegador, o editor mais conhecido é o Overleaf\footnote{\url{https://www.overleaf.com/}}\footnote{Basicamente um Google Docs do \LaTeX{}.}, que, para este modelo, você precisa garantir que esteja sendo usado o compilador XeLaTeX, com o documento aberto, clicando em Menu, no canto superior esquerdo, e selecionando o XeLaTeX como o compilador. E, para utilizar localmente, na sua máquina, existe a opção de utilizar o próprio VS Code como o editor de texto. Essa última opção é bastante vantajosa, por permitir você realizar o controle do seu diretório \textit{git}, se quiser criar, não depender de uma conexão estável de internet\footnote{Eu nunca tive internet decente e, na moral, se você tiver internet horrorosa também, utilize pelo VS Code. Além de que é muito gostosinho de escrever com o auxílio de sugestão de código.} e ter o benefício da usabilidade da interface da IDE e da própria extensão\footnote{Sendo honesto, eu consigo ser muito mais produtivo pelo VS Code do que usando o Overleaf pelo navegador, porém, entendo as vantagens do Overleaf, de permitir compartilhar com outras pessoas e editar simultaneamente. Apesar que, em teoria, a edição simultânea pode ser realizada por causa que existem extensões no VS Code para realizar programação de forma remota e simultânea (VS Code humilha o Overleaf. Venha, junte-se a nós).}.

Para utilizar o \LaTeX{} pelo VS Code, baixe e instale uma instância de \LaTeX{}, como a MiKTeX\footnote{É a que eu utilizo e recomendo. Ela já vem com alguns pacotes úteis instalados por padrão.} e instale a extensão LaTeX Workshop no VS Code. Após isso, ainda é necessário realizar a configuração da extensão: clique \texttt{ctrl+shift+p}, selecione ``\textit{Preferences: Open User Settings (JSON)}'' e copie as configurações que estão no arquivo \texttt{docs/settings.json}, porém, copie com cuidado, selecionando com atenção apenas as configurações relacionadas à extensão e ao \LaTeX{}. A própria extensão LaTeX Workshop possui uma \textit{wiki} que pode servir de auxílio para os ajustes.

\section{Oneside e Twoside}
Uma das configurações possíveis de se realizar no documento é a criação de um documento que utilize os dois versos de uma folha, como um livro. Apesar do \LaTeX{} ter a limitação de não indicar à impressora que ela deve imprimir nos dois lados da folha \cite{site:latex-impressora}, o modelo foi configurado para poder ser utilizado dessa forma, mesmo que não seja algo comum e seja, aparentemente, bastante opcional\footnote{Não garanto que funcione 100\%, não me cobrem se der algo errado.}. E essa configuração ocorre ao optar, no \texttt{documentclass} do documento, entre \texttt{oneside} e \texttt{twoside}.

O Guia de Normalização recomenda, para o caso de utilizar o verso das folhas, que os elementos pré-textuais e pós-textuais sempre iniciem no anverso da folha (a parte da frente). Aparentemente, nesse caso, realmente é apenas uma recomendação, segue quem quer e só se utilizar a opção \texttt{twoside}. Para seguir essa recomendação, foi configurado o uso do \verb|\cleardoublepage|\footnote{O modelo detecta qual configuração está sendo usada para ativar o comando} no fim de cada elemento pré-textual\footnote{Existe uma exceção envolvendo a folha de rosto, que deve ter a ficha catalográfica ou de identificação de obra no seu verso}, e textual, no qual o usuário deverá adicionar o comando em cada um de seus capítulos. Porém, existe um \textit{bug} identificado, que está ocorrendo no último capítulo, que o comando está criando uma página em branco a mais. Nesse caso, apenas deixe comentado o comando, pois as referências, ao se utilizar como capítulo, cria-se sempre em uma nova página\footnote{E se ocorrer algum outro tipo de bug assim, só comente o comando que cria essa página em branco, pois, ele não é tão essencial. Terminar um capítulo \texttt{clearpage} ou \texttt{cleardoublepage} é mais uma boa prática do que necessidade, já que o capítulo sempre inicia em uma nova página.}.


%%%%%%%%%%%%%%%%%%%%%%%%%%%%%%%%%%
% Isso aqui é só um elemento estético, para separar o meu alerta. Por favor, não utilize no seu trabalho algo assim, a ABNT iria chorar se o fizesse.
\vspace{\baselineskip}
\begin{Center}
    \rule{12cm}{1pt}
\end{Center}

\textbf{ATENÇÃO:} É muito comum que um documento no \LaTeX{}, durante a compilação, dê algum erro e ele, na verdade, seja originado por causa de algum caractere diferente do habitual. Alguns caracteres especiais, como o \_, que deve ser utilizado através do comando \verb|\textunderscore ou \_|\footnote{Na Tabela \ref{tabela:exemplo-caracteres} mostra-se um conjunto de caracteres especiais, alguns deles podem levar a erros se digitados diretamente no texto.}, levam a erros na compilação, e podem estar contidos em alguma de suas referências utilizadas, e geralmente estão. Por isso recomenda-se bastante atenção no registro das referências bibliográficas, com exceção de URLs, pois, muitas vezes, esses erros ocorrem de forma silenciosa, sem dar um alerta do que pode estar causando ele (principalmente se você estiver utilizando o método de compilação \textbf{XeLaTeX $\times$ Biber $\times$ XeLaTeX $\times$ XeLaTeX}, que basicamente oculta mensagens de erro quando eles estão acontecendo na bibliografia, por causa das compilações seguintes).
\endgroup

% Seção ilustrativa sobre seções e subseções
\begingroup
  \chapter{Exemplo de seção primária} 
  \lipsum[1]

  \section{Exemplo de seção secundária}
  \lipsum[2-4]

  \subsection{Exemplo de seção terciária}
  \lipsum[5]

  \subsubsection{Exemplo de seção quaternária}
  \lipsum[6]

  \paragraph{Exemplo de seção quinária} % Utiliza-se paragraph para o nível hierarquico inferior à subsubsection. Contudo, foi necessário customizá-lo para apresentar corretamente como um nível de seção
  \lipsum[7]

  \paragraph{Exemplo de título com indicação numérica que, ao ocupar mais de uma linha, deve ser, a partir da segunda linha, alinhado abaixo da primeira letra da primeira palavra do título}
  \lipsum[8]
\endgroup

\begingroup
  \chapter{Exemplos de elementos comuns}
\section{Alíneas e subalíneas}
  \lipsum[1-1]
  \begin{alinea}
    \item a matéria da alínea começa por letra minúscula, exceto quando se tratar de     substantivos próprios, e termina em ponto e vírgula, com exceção da última, que termina em ponto final; 
    \item o trecho final da seção correspondente, anterior às alíneas, termina em
    dois pontos;
    \item asdasdsada: 
    \begin{subalinea}
      \item a matéria da alínea começa por letra minúscula, exceto quando se tratar de     substantivos próprios, e termina em ponto e vírgula, com exceção da última, que termina em ponto final; 
      \item o trecho final da seção correspondente, anterior às alíneas, termina em
      dois pontos;
    \end{subalinea}
    \item sdsdsdaa
  \end{alinea}

  \subsection{Outros exemplos de listas numeradas}
    % \begin{leiParagrafo}
    %   \item a matéria da alínea começa por letra minúscula, exceto quando se tratar de     substantivos próprios, e termina em ponto e vírgula, com exceção da última, que termina em ponto final; 
    %   \item o trecho final da seção correspondente, anterior às alíneas, termina em
    %   dois pontos;
    %   \item asdasdsada
    %   \item sdsdsdaa
    % \end{leiParagrafo}

\section{Figuras, tabelas e quadros}
  --- Falar que figure, table e quadro são apenas invólucros, que qualquer coisa pode ser adicionada dentro, que apenas indica que o elemento possuirá seu respectivo tratamento diferenciado na hora de listar e apresentar como tal
  
\subsection{Figuras}
  --- Falar os formatos de imagem suportados
  \begin{figure}[H]
    \centering\singlespacing
    \caption{As dimensões da comunicação organizacional}
    \label{figura:marca-iffar}
    % Podem ser utilizados vários argumentos para adicionar a figura em seu devido tamanho, como: scale; width; ou height. No caso de scale, a escala, adiciona o valor em decimal representando a porcentagem da escala, já width ou height utiliza-se uma unidade de medida, como cm, px ou pt.
    \includegraphics[scale=0.35]{marca-iffar.png}\\
    \footnotesize
    Fonte: \textcite{iffar-identidade-visual-2021}.
  \end{figure}
  
\subsection{Tabelas}
\begin{table}[H]
  \centering\singlespacing

  \caption{Posições do IFFar no \textit{ranking} do painel de dados da LAI}
  \label{tabela:cgu-iffar}
  \begin{tabularx}{12cm}
    {X r r} % Definição das colunas
    \hline
    \multicolumn{1}{c}{\textbf{Categoria}}    &   
    \multicolumn{1}{c}{\textbf{Posição}}  & 
    \multicolumn{1}{c}{\textbf{Informação adicional}} \\
    \hline

    Transparência ativa &
    252\textordmasculine{} &
    24,49\% em cumprimento \\
    %\hline

    Pedidos recebidos &
    192\textordmasculine{} &
    807 pedidos \\
    %\hline

    Tempo médio de resposta para pedidos de informação       &   
    154\textordmasculine{} &
    14,91 dias \\
    \hline

  \end{tabularx}

\hspace{\fill}

\footnotesize

Fonte: Gaba guba.
\end{table}

\subsection{Quadros}
--- Falar que é uma tabela com bordinha nos lados
\begin{quadro}[H]
  \centering\singlespacing

  \caption{Posições do IFFar no \textit{ranking} do painel de dados da LAI}
  \label{quadro:cgu-iffar}
  \begin{tabularx}{12cm}
    {|X|r|r|} % Definição das colunas
    \hline
    \multicolumn{1}{|c|}{\textbf{Categoria}}    &   
    \multicolumn{1}{c}{\textbf{Posição}}  & 
    \multicolumn{1}{|c|}{\textbf{Informação adicional}} \\
    \hline

    Transparência ativa &
    252\textordmasculine{} &
    24,49\% em cumprimento \\
    %\hline

    Pedidos recebidos &
    192\textordmasculine{} &
    807 pedidos \\
    %\hline

    Tempo médio de resposta para pedidos de informação       &   
    154\textordmasculine{} &
    14,91 dias \\
    \hline

  \end{tabularx}

\hspace{\fill}

\footnotesize

Fonte: Gaba guba.
\end{quadro}
    

  Conforme o quadro no Apêndice \ref{apendice:teste} e o Anexo \ref{anexo:teste}. Citando \textcite{brasil-lei-11.892-2008}
  
  Testesddsdfdfs

\endgroup

%% Aqui são adicionados as referências, apêndices e anexos. É essencial que eles estejam dentro desse ambiente, para formatar corretamente o título da seção deles
\begin{SecoesNaoNumeradas}
%% Aqui é impressa a lista de referências bibliográficas
\begingroup
  \begin{FlushLeft} % As referências devem apresentar alinhamento à esquerda
    \begin{singlespace} % Referências usam espaço simples
      \setlength\bibitemsep{\baselineskip} % Defino que o espaço para separar cada referência é o da altura de uma linha, conforme a ABNT exige
      % heading=bibintoc imprime o nome da seção das referências sem numeração e adiciona no sumário
      % title=Referências modifica o nome da seção, que normalmente seria "Bibliografia", por conta que é o nome utilizado pelos documentclass book e report
      \printbibliography[heading=bibintoc, title=Referências] % ITEM OBRIGATÓRIO
    \end{singlespace}
  \end{FlushLeft}
\endgroup

% Glossário não tem muito mistério, é só uma lista de termos e suas explicações
\chapter*{Glossário}
\noindent Exemplo de termo: explicação.

\noindent Exemplo de outro termo: Lorem ipsum dolor sit amet, consectetur adipiscing elit.

%% Depois das referências, apresentam-se os apêndices e, depois, os anexos
\begingroup
\apendice{Um exemplo de apêndice}\label{apendice:teste}
Alíneas\index{alínea}
\lipsum[6-7]
\endgroup

\begingroup 
\anexo{Um exemplo de anexo}\label{anexo:teste}
\lipsum[6-7]
\endgroup

% \chapter*{Índice}
\begin{indice}
    \item Teste
    \item Exemplo
\end{indice}
\printindex % Comente esta linha se não for utilizar índice

\end{SecoesNaoNumeradas}

\end{document}